This explanation for the non-existence of odd-parity pulsations can be restated in
more physical terms as follows: Odd-parity gravitational waves are ``transverse'' waves
in the sense that they only couple to stars which can support anisotropic stresses.
(Anisotropic stresses are characterized by tensors, whose spherical harmonics can have
odd parity.) Consequently, by idealizing our equilibrium configurations as made of a
perfect fluid (purely isotropic stress, $-p$), we rule out, ab initio, the possibility of
odd-parity pulsations and of the generation of odd-parity gravitational waves.

Because odd-parity pulsations are non-existent, we shall confine our analysis to
even-parity motions throughout the remainder of this paper.

\subsection*{(d) Even-Parity Motions}

The equations of motion which govern the time evolution of even-parity motions are
derived from Einstein's field equations in Appendix~C\@. For a given harmonic
$(l,\,M=0,\,s=(-1)^{l})$ these equations of motion are a set of coupled equations in the
fluid displacement functions $V(r,t)$, $W(r,t)$ and for the metric perturbation functions
$H_0(r,t)$, $H_1(r,t)$, $H_2(r,t)$, $K(r,t)$ (cf.\ eqs.\ [3]). Initial-value equations, which
determine the time evolution of $K$, $W$, $V$ from initial data, have been specified; and
propagation equations, which determine the time evolution of $K$, $W$, $V$. For $l\geq 2$
the \emph{initial-value equations} are (cf.\ Appendix~C)

\begin{equation}
\begin{split}
H_0' + r^{-1}e^{\lambda}\bigl[l(l+1)/2 + 1 + 4\pi r^2(\rho-p)\bigr]H_0 \\
\quad = rK'' + e^{\lambda}(3-5m/r-4\pi r^2 p)K'
  - r^{-1}e^{\lambda}\bigl[l(l+1)/2 - 1 \\
\qquad + 8\pi r^2(\rho+p)\bigr]K
  + 8\pi r^{-1}(\rho+p)e^{\lambda/2}W' \\
\qquad + 8\pi r^{-1}\rho'\,e^{\lambda/2}W
  + 8\pi l(l+1)r^{-1}(\rho+p)e^{\nu}V \;,
\end{split}
\tag{8a}
\end{equation}

\begin{equation}
\begin{split}
l(l+1)H_1 = 2r^2 K'_{,t} - 2re^{\lambda}(-1+3m/r+4\pi r^2 p)K_{,t}
  - 2rH_{0,t} \\
  + 16\pi(\rho+p)e^{\lambda/2}W_{,t} \;,
\end{split}
\tag{8b}
\end{equation}

\begin{equation}
H_{1,t} = e^{\nu}H_0' + r^{-1}e^{\lambda+\nu}(2m/r+8\pi r^2 p)H_0 - e^{\nu}K' \;,
\tag{8c}
\end{equation}

\begin{equation}
H_2 = H_0 \;,
\tag{8d}
\end{equation}

and the time derivatives of equations (8a) and (8d). For $l\geq 2$ the \emph{propagation
equations} are

\begin{equation}
\begin{split}
K_{,tt} - e^{-\lambda}K'' + 2r^{-1}e^{-\lambda}\bigl[-1+m/r+2\pi r^2(\rho-p)\bigr]K' \\
+ r^{-2}e^{\nu}\bigl[l(l+1)-2+8\pi r^{-2}(\rho+p-\gamma p)\bigr]K
  + r^{-2}e^{\nu}\bigl[2e^{-\lambda}-4\pi r^2(\rho+p+\gamma p)\bigr]H_0 \\
- 8\pi r^{-2}(\rho+p-\gamma p)e^{\nu-\lambda/2}W'
  - 8\pi r^{-2}(\rho'-p')e^{\nu-\lambda/2}W \\
- 8\pi l(l+1)r^{-2}(\rho+p-\gamma p)e^{\nu}V = 0 \;,
\end{split}
\tag{9a}
\end{equation}

\begin{equation}
\begin{split}
r^{-2}(\rho+p)e^{(\lambda-\nu)/2}W_{,tt}
  = \bigl\{\gamma p e^{\nu/2}\bigl[r^{-2}e^{-\lambda/2}W'
    + l(l+1)r^{-2}V - \tfrac{1}{2}H_0 - K\bigr]\bigr\}' \\
+ \tfrac{1}{2}(\rho+p)e^{\nu/2}(r^{-2}e^{-\lambda/2}\nu')'W
  - l(l+1)r^{-2}(\rho+p)(e^{\nu/2})'V \\
- \tfrac{1}{2}(\rho+p)e^{-\nu}(H_0 e^{3\nu/2})'
  + (\rho+p)(Ke^{\nu/2})' = 0 \;,
\end{split}
\tag{9b}
\end{equation}

\begin{equation}
\begin{split}
e^{-\nu}(\rho+p)V_{,tt} + l(l+1)r^{-2}\gamma p V
  + r^{-2}\gamma p\,e^{-\lambda/2}W' + r^{-2}p'\,e^{-\lambda/2}W \\
- \tfrac{1}{2}(p+\rho+\gamma p)H_0 - \gamma p K = 0 \;.
\end{split}
\tag{9c}
\end{equation}
